\section{Methodology}
\label{sec:methodology}

We conduct three experiments on \gpttwo~\citep{radford2019language} (124M parameters, 12 layers, 768 dimensions), accessing internal activations via \transformerlens~\citep{nanda2022transformerlens} and using pre-trained SAEs (24,576 features per layer) from the GPT-2 Small SAE suite.
All experiments run on a single NVIDIA RTX A6000 (48GB) in approximately 3 minutes total.

\subsection{Dataset Construction}
\label{sec:dataset}

We construct a dataset of 21 compound concepts spanning the compositionality spectrum, with ratings from \emph{very\_low} (fully idiomatic, \eg ``red herring'') to \emph{very\_high} (fully transparent, \eg ``kitchen chair'').
\tabref{tab:compounds} shows representative examples.

\begin{table}[t]
    \centering
    \begin{tabular}{@{}llllc@{}}
        \toprule
        \textbf{Compound} & \textbf{Modifier} & \textbf{Head} & \textbf{Category} & \textbf{Comp.} \\
        \midrule
        washing machine & washing & machine & appliance & medium \\
        hot dog & hot & dog & food & low \\
        red herring & red & herring & idiom & very\_low \\
        kitchen chair & kitchen & chair & furniture & very\_high \\
        swimming pool & swimming & pool & structure & high \\
        guinea pig & guinea & pig & animal & low \\
        dark horse & dark & horse & idiom & very\_low \\
        coffee machine & coffee & machine & appliance & high \\
        school bus & school & bus & vehicle & high \\
        ice cream & ice & cream & food & medium \\
        \bottomrule
    \end{tabular}
    \caption{Representative compounds from our dataset with compositionality ratings. Full dataset contains 21 compounds spanning appliances, food, vehicles, idioms, furniture, and clothing categories.}
    \label{tab:compounds}
\end{table}

Each compound is processed in the minimal context ``The [compound]'' to isolate the compound representation from richer contextual effects.
For component baselines, we process ``The [modifier]'' and ``The [head]'' separately.

\para{Tokenization considerations.}
GPT-2's \BPE tokenizer handles these compounds variably: ``washing machine'' splits cleanly into two tokens (\texttt{washing}, \texttt{\textvisiblespace machine}), but ``coffee machine'' becomes three tokens (\texttt{co}, \texttt{ffee}, \texttt{\textvisiblespace machine}) and ``red herring'' splits the head noun (\texttt{red}, \texttt{\textvisiblespace her}, \texttt{ring}).
For multi-token components, the model must already perform compositional processing just to represent the component word, adding complexity to the analysis.
We extract representations at the last token position for all conditions.

\subsection{Experiment 1: Residual Stream Cosine Similarity}
\label{sec:exp1}

For each compound concept, we measure how the compound representation relates to its components in the residual stream at every layer.
Specifically, at each layer $\ell$, we compute three cosine similarities:
\begin{align}
    \cossim_{\text{head}} &= \frac{\vh^{(\ell)}_{\compound} \cdot \vh^{(\ell)}_{\text{head}}}{\|\vh^{(\ell)}_{\compound}\| \|\vh^{(\ell)}_{\text{head}}\|} \label{eq:cos_head} \\
    \cossim_{\text{add}} &= \frac{\vh^{(\ell)}_{\compound} \cdot \bar{\vh}^{(\ell)}_{\text{avg}}}{\|\vh^{(\ell)}_{\compound}\| \|\bar{\vh}^{(\ell)}_{\text{avg}}\|} \label{eq:cos_add} \\
    \cossim_{\text{mod}} &= \frac{\vh^{(\ell)}_{\compound} \cdot \vh^{(\ell)}_{\text{mod}}}{\|\vh^{(\ell)}_{\compound}\| \|\vh^{(\ell)}_{\text{mod}}\|} \label{eq:cos_mod}
\end{align}
where $\vh^{(\ell)}_{\compound}$ is the hidden state at the last token of the compound, $\vh^{(\ell)}_{\text{head}}$ and $\vh^{(\ell)}_{\text{mod}}$ are the hidden states for the head noun and modifier processed in isolation, and $\bar{\vh}^{(\ell)}_{\text{avg}} = \frac{1}{2}(\vh^{(\ell)}_{\text{mod}} + \vh^{(\ell)}_{\text{head}})$ is the additive composition baseline.

\subsection{Experiment 2: SAE Feature Decomposition}
\label{sec:exp2}

We use pre-trained SAEs at layers 1, 6, and 11 to decompose residual stream activations into sparse feature sets.
Each SAE has 24,576 features per layer.
For each compound and its components, we:
\begin{enumerate}[leftmargin=*,itemsep=0pt,topsep=0pt]
    \item Encode the residual stream activation through the SAE encoder.
    \item Identify active features (activation $> 0$).
    \item Compute the fraction of compound features that are \emph{unique to the compound context}---not shared with either the head or modifier processed alone.
    \item Compute Jaccard similarity $J = |F_{\compound} \cap F_{\text{head}}| / |F_{\compound} \cup F_{\text{head}}|$ between compound and head feature sets.
\end{enumerate}

This decomposition tests whether the subtle modifications that cosine similarity misses correspond to genuine differences in feature activation patterns.

\subsection{Experiment 3: Next-Token Prediction Priming}
\label{sec:exp3}

For each compound ``modifier head,'' we measure how strongly the modifier primes the head noun in the model's next-token predictions:
\begin{equation}
    \text{Priming Ratio} = \frac{P(\text{head} \mid \text{``The modifier''})}{P(\text{head} \mid \text{``The''})}
    \label{eq:priming}
\end{equation}
where the numerator is the probability of the head noun after seeing the modifier in context, and the denominator is its baseline probability after seeing only ``The.''
We also record the rank of the head noun among all 50,257 vocabulary predictions.
This experiment directly tests whether the model ``stores'' compound concepts primarily through transition probabilities rather than dedicated feature directions.
